% ECONOMICS_PROBLEM: {"id": "C73-28", "title": "Nash equilibrium in two-player concurrent reachability games", "jel_code": "C73", "source_id": "OP-0158"}
% !TeX program = xelatex
\documentclass[11pt,a4paper]{article}
\usepackage[margin=23mm]{geometry}
\usepackage{fontspec}
\setmainfont{Latin Modern Roman}
\usepackage{amsmath,amssymb,mathtools}
\usepackage{xcolor,enumitem,booktabs,longtable,array,xurl,hyperref}
\definecolor{muted}{HTML}{53636C}
\hypersetup{colorlinks=true,urlcolor=blue,linkcolor=blue}
\setlength{\parindent}{0pt}
\setlength{\parskip}{5pt}
\newcommand{\R}{\mathbb R}
\newcommand{\E}{\mathbb E}
\newcommand{\N}{\mathbb N}
\newcommand{\ind}{\mathbf 1}
\newcommand{\Var}{\operatorname{Var}}
\newcommand{\Cov}{\operatorname{Cov}}
\newcommand{\supp}{\operatorname{supp}}
\DeclareMathOperator*{\argmin}{arg\,min}
\DeclareMathOperator*{\argmax}{arg\,max}
\newcommand{\entrymeta}[3]{{\sffamily\small\color{muted}\textbf{\texttt{#1}}\quad|\quad #2\par #3\par}}
\newcommand{\Problemref}[1]{\texttt{#1}}
\newcommand{\JELref}[2]{\texttt{#2} (JEL \texttt{#1})}
\begin{document}
\section*{Nash equilibrium in two-player concurrent reachability games}
\textbf{Problem ID:} \texttt{C73-28}\quad \textbf{JEL:} \texttt{C73}\par
% BEGIN ECONOMICS STATEMENT
\label{problem:OP-0158}
% GAME_THEORY_PROBLEM: {"local_id": "GT-028", "original_number": 28, "kind": "P", "entry_kind": "statement_only", "published": false, "review_required": true, "category": "game_theory"}
\label{game:GT-028}
{\sffamily\small\color{muted}
{\sffamily\small\color{muted}\textbf{GT-028} | Stochastic, repeated, and dynamic games\par P | Source-reported exact existence question\par}
\par}
\subsection*{Definitions and mathematical statement}
Let $S$ be finite, $s_0\in S$, $A_i(s)$ nonempty and finite for $i\in\{1,2\}$, and $P(s'\mid s,a_1,a_2)$ a transition kernel. Both players choose actions simultaneously at each stage. They observe the public state history. A strategy is a map $\sigma_i:S^+\to\Delta(A_i(s))$, where $s$ is the last state of the input history; private action choices are not additional inputs to this source-convention strategy class. For target sets $F_i\subseteq S$, put
\[
 u_i(\sigma)=\Pr_{s_0,\sigma}(\exists t\geq0:\ s_t\in F_i).
\]
The exact Nash-equilibrium question is
\[
 \forall G\ \forall s_0\quad\exists(\sigma_1,\sigma_2)\quad
 \forall i\in\{1,2\}\ \forall\tau_i:S^+\to\Delta(A_i),
 \qquad u_i(\tau_i,\sigma_{-i})\leq u_i(\sigma).
\]
The displayed shorthand for $\tau_i$ always requires actions feasible at the last state. The induced probability law is obtained by independently sampling the two prescribed action distributions conditional on the observed state history, then applying $P$.
\subsection*{Short English statement}
Does every finite two-player concurrent game with reachability payoffs possess an exact Nash equilibrium?
\subsection*{Variants and scope boundaries}
Targets need not be absorbing. Exact existence, approximate existence, and stationary existence differ. Observing joint actions in addition to states changes the strategy domain; so does restricting strategies to finite transducers. For complexity problems, rational versus dyadic transitions and their encodings must be fixed separately.
\subsection*{Source relationship and status}
[S1], Section 1, reports this question as open and explains that an earlier proposed counterexample was flawed. Section 2 defines the state-history monitoring convention used here. Its finite-horizon equilibrium verification results do not settle infinite-horizon existence.
\subsection*{References}
\begingroup\footnotesize\raggedright
\textbf{[S1]} Senthil Rajasekaran and Moshe Y. Vardi, \emph{Verifying Equilibria in Finite-Horizon Probabilistic Concurrent Game Systems}, arXiv:2503.14690v7 (2026 revision of a 2025 preprint). Locators: Section 1, paragraphs on infinite-horizon realizability and the flawed counterexample; Section 2, strategy and Nash-equilibrium definitions. \url{https://arxiv.org/html/2503.14690v7}
\par\endgroup

\subsection*{Common conventions: Source mathematical conventions}

\textbf{Finite games.} For a finite set $A$, $\Delta(A)$ is its probability
simplex. Players choose independent mixed strategies in Nash equilibrium;
correlation is allowed only when explicitly part of the solution concept.
In a game with payoff functions $u_i$ and strategy profile $x$, an additive
$\varepsilon$ Nash equilibrium satisfies
\[
 u_i(x)\geq u_i(x_i',x_{-i})-\varepsilon
 \quad\text{for every player $i$ and every admissible deviation $x_i'$}.
\]
For numerical approximation, payoffs are normalized as locally specified.
An $\varepsilon$ payoff residual is distinct from distance at most
$\varepsilon$ to the set of exact equilibria. Point convergence, convergence
of empirical play, and convergence in distance to an equilibrium set are
also distinct.

\textbf{Computational representations.} Unless stated otherwise, a complexity
question takes rational data with binary encoding; polynomial time is measured
in total bit length. A fixed-accuracy polynomial algorithm is not necessarily
polynomial in $1/\varepsilon$, and neither is necessarily polynomial in
$\log(1/\varepsilon)$. Succinct games and value, demand, or sampling oracles
must declare their interfaces. Exact real solutions require an explicit
representation; an unspecified list of arbitrary real numbers is not a
finite computational output. A claimed algorithmic barrier is meaningful
only for its specified model and reductions.

\textbf{Dynamic games.} Histories, information, observation, and allowable
randomization are part of the model. A stationary strategy depends only on
the current state; a positional strategy in a deterministic graph game
chooses one action at each controlled vertex. Nash equilibrium at an initial
state and equilibrium after every history are different requirements.
An expectation of a pathwise $\liminf$ payoff, a $\liminf$ of expected averages,
a limiting expected average, and a uniform finite-horizon equilibrium are
different criteria. None is silently substituted for another.

\textbf{Allocation and mechanisms.} A complete allocation assigns every item
exactly once unless otherwise stated. Zero-valued goods, negative values,
capacity constraints, feasibility, lotteries, and copies of goods affect
fairness definitions and are specified locally. Dominant-strategy incentive
compatibility, Bayesian incentive compatibility, universal truthfulness,
and truthfulness in expectation impose different inequalities. Whenever
an objective ratio has zero denominator, its convention is stated or those
instances are excluded.

\textbf{Infinite-dimensional models.} Expectations, integrals, stochastic
equations, and optimization objectives require their local regularity,
integrability, filtrations, and admissible domains. A backward--forward
mean-field system is a boundary-value problem; it is not an initial-value
semiflow without an additional construction. Long-time, large-population,
and vanishing-noise limits require an order of limits and a convergence
topology. If a minimum need not be attained, the objective is its infimum
or supremum and the approximation requirement is stated separately.
% END ECONOMICS STATEMENT
\end{document}
